\documentclass[11pt]{article}

\usepackage[margin=1in]{geometry}
\usepackage[T1]{fontenc}
\usepackage[utf8]{inputenc}
\usepackage{microtype}
\usepackage{hyperref}
\usepackage{enumitem}
\usepackage{booktabs}
\usepackage{amsmath}
\usepackage{amssymb}
\usepackage{array}
\usepackage{xcolor}
\setlength{\emergencystretch}{2em}

\hypersetup{
  colorlinks=true,
  linkcolor=black,
  citecolor=black,
  urlcolor=black,
  pdftitle={From Episodes to Primitives},
  pdfauthor={Keren Dong}
}

\newcommand{\kfd}{\textsc{KFD}}
\newcommand{\episode}{\textsc{Episode}}

\title{From Episodes to Primitives:\\
Causal Historical Assets for Evidence-Governed Ontology Revision}
\author{Keren Dong\\Kungfu Origin Technology Limited\\\href{mailto:keren.dong@kungfu.link}{keren.dong@kungfu.link}}
\date{Alpha draft: July 2026}

\begin{document}

\maketitle

\begin{abstract}
Agent systems usually preserve histories to debug prior actions or retrieve
prior answers. This paper studies a stronger possibility: a well-formed causal
\episode{} corpus may preserve the option to discover that the ontology under
which earlier questions were asked was incomplete. We define an Episode as a
perspective-bound, provenance-carrying record of observations, choices,
actions, dependencies, consequences, interventions, omissions, and adjacent
fact cuts. We then specify a staged mechanism derived from Kung Fu Decisions
(KFD): KFD-4 replays causal histories from declared perspectives; KFD-5
separates candidate genesis from qualification; draft KFD-6 compares plural
generation methods against fixed-ontology and no-new-primitive baselines,
held-out evidence, and separated promotion authority. The resulting asset is
not an event pile, transcript, embedding index, or automatic proof. Its future
value depends on integrity, perspective, causal and responsibility boundaries,
portability, retention, privacy, consent, and independent verification. We
ground the design in one immutable July 2026 KFD evidence cut containing three
live cases and six candidate tracks: three accepted only within software-domain
profiles and three provisional. This is early feasibility evidence, not a
conforming KFD-6 experiment. We contribute an artifact taxonomy, an
evidence-governed discovery loop, a governance analysis of problem-definition
power, and a falsifiable prospective evaluation protocol. The central claim is
conditional: governed causal history can preserve future ontology-revision
capacity; storage volume, narrative coherence, or generator confidence cannot.
\end{abstract}

\section{Introduction}
\label{sec:introduction}

Most computational histories are collected to answer questions already known
at capture time: which call failed, which state changed, which action produced
a reward, or which artifact was released. The schema names the objects in
advance. Later work searches, aggregates, or predicts within that object world.

Real work creates a harder failure mode. A team can preserve every field it
knows and still lose the relation that a future participant needs, because the
relation did not yet have a name. Permission can be mistaken for accepted
responsibility. An executed action can be mistaken for a settled obligation.
A collection of verified components can be mistaken for the one project state
a successor is authorized to continue. The missing information is not always a
missing value. It may be a missing object boundary.

This paper asks whether historical infrastructure can preserve a capacity to
repair that kind of omission. The proposal builds on three Kung Fu Decisions:
KFD-4 binds views to declared perspectives and permits source-preserving replay
and transformation; KFD-5 separates candidate genesis from fact-bound
qualification; draft KFD-6 proposes a bounded loop that compares generation
methods over causal experience without allowing a generator to certify or
promote itself \cite{kfd4Decision2026,kfd5Decision2026,kfd6Decision2026}.

The central thesis is conditional:

\begin{quote}
A well-formed causal Episode corpus can preserve more than prior answers. It
can preserve a governed option to replay decisions, expose unnamed burdens,
qualify candidate Primitives, and revise the ontology under which future
questions are formed.
\end{quote}

The word \emph{option} matters. A corpus does not discover merely by existing.
A generated label is not a Primitive. A plausible retrospective story is not
causal evidence. The proposed value appears only when later participants can
identify an immutable evidence cut, reconstruct perspectives with visible
loss, compare plural methods and conservative baselines, reject attractive
false candidates, and keep promotion authority separate from generation.

The contribution is distinct from two adjacent KFD papers. The Foundation
paper develops facts, trust, cooperation, and boundary primitives as an ordered
systems model \cite{kfdFoundationPaper2026}. The Observer paper develops
declared timelines and perspective-preserving replay
\cite{kfdObserverPaper2026}. This paper treats replayable Episodes as a
long-lived input to ontology revision and focuses on the KFD-4 to KFD-5 to
KFD-6 transition.

We make five contributions:

\begin{enumerate}[leftmargin=*]
  \item We define an \episode{} historical asset and distinguish it from a
        transcript, event log, provenance graph, embedding index, and endpoint
        state.
  \item We specify a staged discovery mechanism that keeps perspective replay,
        candidate genesis, qualification, held-out evaluation, and promotion
        authority independently inspectable.
  \item We explain why control of historical representation affects the power
        to define future questions, while deriving stewardship obligations
        rather than a command to retain everything.
  \item We freeze one implementation-backed KFD evidence cut and grade what it
        supports: three live cases and six tracks, with three bounded
        software-domain acceptances and three provisional candidates.
  \item We propose prospective ablation, shadow-loop, transfer, and governance
        evaluations capable of falsifying the stronger thesis.
\end{enumerate}

The paper does not claim autonomous Primitive discovery, universal ontology,
historical novelty of current KFD candidates, or a quantified return on
Episode retention. It claims that a mechanism can be specified, its early
feasibility can be separated from its missing gates, and the next experiment
can be made answerable to evidence.

\section{Problem Statement and Research Questions}
\label{sec:problem}

Let an operational system at time $n$ act with ontology $O_n$. Its capture
schema determines which entities, relations, state changes, and responsibilities
are directly representable. A fixed-ontology history $H(O_n)$ can be perfectly
internally consistent while remaining unable to express a distinction $x$
that becomes load-bearing later.

An ontology revision is justified only when the new distinction changes a
required decision and cannot be recovered from the old representation at lower
total burden. The problem is therefore not to maximize labels. It is to detect
when repeated reconstruction, hidden coordination, divergent consequences, or
causal variables indicate that $O_n$ is missing an independently addressable
responsibility.

We study four research questions:

\begin{description}[leftmargin=2.1em,style=nextline]
  \item[RQ1: Representation.] Which properties distinguish an Episode corpus
  that preserves future ontology-revision capacity from anonymous traces or
  endpoint facts?
  \item[RQ2: Mechanism.] Can declared perspective replay improve candidate
  genesis while KFD-5 qualification prevents narrative or method artifacts
  from becoming Primitives?
  \item[RQ3: Evaluation.] Under shared evidence and resource budgets, does an
  Episode-based loop outperform fixed-ontology and no-new-Primitive baselines
  on held-out reconstruction, transfer, and false-candidate rejection?
  \item[RQ4: Governance.] Which integrity, retention, deletion, consent,
  portability, access, and promotion boundaries are necessary when a corpus
  may influence not only answers but future problem definitions?
\end{description}

These questions deliberately separate technical possibility from legitimate
authority. A system may generate a useful candidate and still lack the right
to retain the underlying experience, represent a participant, or change an
adopted ontology.

\section{Episode Historical Assets}
\label{sec:episode-assets}

\subsection{From recorded events to causal experience}

Event order is necessary but insufficient. Distributed histories already show
that order and causality cannot be inferred from a single wall clock
\cite{lamport1978time}. Provenance standards distinguish entities, activities,
agents, derivation, and responsibility \cite{w3cProv2013}. KFD adds a
participant and ontology boundary: what did an observer accept, which action
was attempted under which authority, what consequence was borne, what evidence
was missing, and which later interpretation remains a claim?

We model an Episode $e$ as:

\begin{equation}
e = (i, P, F^{-}, A, D, C, I, F^{+}, R, G, S),
\end{equation}

where $i$ is stable identity; $P$ is declared observer and perspective;
$F^{-}$ and $F^{+}$ are adjacent admitted Fact cuts; $A$ is the action and
bounded intent; $D$ records dependencies and authority; $C$ records
consequences, costs, failures, retries, and overrides; $I$ records
interventions; $R$ binds evidence and receipts; $G$ declares gaps, redaction,
degradation, and proxy-only consequences; and $S$ declares stewardship policy.

This tuple is a responsibility decomposition, not a required physical layout.
An implementation may use several stores or one content-addressed bundle. It
qualifies as an Episode only when the relevant distinctions remain recoverable
and verifiable.

\subsection{Artifact taxonomy}

Table~\ref{tab:taxonomy} prevents a common category error: more storage does
not imply a stronger historical asset.

\begin{table}[ht]
\centering
\small
\begin{tabular}{p{0.16\linewidth}p{0.31\linewidth}p{0.42\linewidth}}
\toprule
Artifact & Useful capability & Insufficient by itself because \\
\midrule
Transcript & Recover utterances and tool output & Speaker position, action
authority, omitted evidence, consequences, and admission may be implicit. \\
Event log & Query ordered or timestamped events & The event schema may flatten
observer, responsibility, causal dependency, and competing views. \\
Provenance graph & Relate entities, activities, agents, and derivations & It
does not necessarily preserve the ontology under test, qualification gate, or
promotion authority. \\
Endpoint Fact cut & State what was admitted at one boundary & It can omit the
path, alternatives, intervention, failure, and consequence that explain why a
new distinction emerged. \\
Derived index & Retrieve summaries, embeddings, or clusters efficiently & It
is a projection whose loss, source cut, update policy, and non-authority must
remain explicit. \\
Episode asset & Replay a perspective-bound causal occurrence with evidence and
governance & It preserves discovery options but neither generates nor
qualifies a Primitive automatically. \\
\bottomrule
\end{tabular}
\caption{Historical artifact taxonomy. The categories may compose; they are
not interchangeable evidence grades.}
\label{tab:taxonomy}
\end{table}

\subsection{Asset conditions}

A corpus $\mathcal{E}_c$ is an Episode historical asset only if its declared
cut $c$ satisfies six conditions:

\begin{enumerate}[leftmargin=*]
  \item \textbf{Integrity:} stable identities, roots, corruption detection,
        recovery boundaries, and immutable prior cuts.
  \item \textbf{Perspective:} observer position, accepted facts, natural
        objects, consequences, and known gaps remain visible.
  \item \textbf{Causal and responsibility structure:} actions, dependencies,
        authority, outcomes, interventions, and omissions are not flattened
        into anonymous tokens.
  \item \textbf{Portability:} export and import preserve exact roots or declare
        loss; provider migration does not silently redefine identity.
  \item \textbf{Governance:} retention, deletion, consent, access, redaction,
        purpose, and review policies are enforceable and auditable.
  \item \textbf{Independent challenge:} another accountable participant or
        verifier can inspect evidence, degradation, and later qualification.
\end{enumerate}

The corpus may be intentionally incomplete. Completeness claims are stronger
than asset claims and usually unavailable. A useful asset makes omissions
inspectable instead of generating certainty to cover them.

\subsection{Option value without storage determinism}

Let $Q_f$ be a future distribution of admissible questions and $\Pi$ the
policies that permit a participant to inspect a corpus. We define the
\emph{discovery option value} of an Episode cut qualitatively as the avoidable
future burden of reconstructing perspective, causality, and responsibility
when a question requires concepts absent from $O_n$:

\begin{equation}
V_{\mathrm{option}}(\mathcal{E}_c) =
\mathbb{E}_{q \sim Q_f}
\left[ B(q \mid H(O_n)) - B(q \mid \mathcal{E}_c, \Pi) \right].
\end{equation}

$B$ may include reconstruction time, coordination load, missing-evidence risk,
or inability to test an ontology revision. The expression is not a monetary
valuation and does not imply $V_{\mathrm{option}}>0$. Poor capture, restricted
access, privacy harm, corrupted roots, or ontology lock-in can make retention
neutral or harmful. Storage volume is absent from the definition because
volume does not establish recoverability or legitimacy.

\section{The KFD-4/5/6 Discovery Mechanism}
\label{sec:mechanism}

The mechanism is a sequence of gates, not one model call. Each stage preserves
a different authority boundary.

\subsection{KFD-4: replay changes the observation horizon}

For admitted facts $F$, a perspective $P$ defines a projection
$\pi_P(F)$. A transformation $T_{PQ}$ reconstructs a new declared view without
claiming to become the source observer:

\begin{equation}
T_{PQ}: (P, \pi_P(\gamma)) \rightarrow
(Q, \pi_Q(\gamma), J, L),
\end{equation}

where $\gamma$ is a causal path, $J$ is the set of preserved invariants, and
$L$ is declared loss. Source identity, Fact and evidence cuts, known causal
relations, receipts, and authority boundaries may belong to $J$. Relevance,
ordering, natural object boundaries, and available action may change.

Contrastive replay keeps at least two source views distinct. Its purpose is not
to vote on one universal view. It asks which burden becomes natural for the
participant who bears the next consequence. A user may expose an object hidden
from an implementation view; a successor agent may expose a settlement object
hidden from the coordinator who assembled its components.

KFD-4 is one genesis method. Direct situated judgment, anomaly search,
reconstruction analysis, causal-variable discovery, and structural compression
remain first class. This method plurality prevents a successful perspective
story from becoming a ranking theorem.

\subsection{KFD-5: generation is not qualification}

Genesis binds the perspective, current ontology, evidence cut, method,
observation, candidate, and claim boundary. Qualification then asks whether the
candidate deserves independent identity, boundary, authority, lifecycle, or
operations.

For candidate $x$, define $\operatorname{remove}_x(h)$ as a history that erases
the semantic distinction $x$ while preserving other observable facts. The
candidate is non-redundant for decision $d$ when two valid histories exist such
that:

\begin{equation}
\operatorname{remove}_x(h_1)=\operatorname{remove}_x(h_2)
\quad \land \quad
D_d(h_1) \ne D_d(h_2).
\end{equation}

This deletion witness does not privilege a name or physical representation. A
rival model passes if it preserves equivalent information with lower total
burden. KFD-5 therefore requires alternatives, minimum closure, deletion and
fuse tests, falsifiers, dogfood under load, residual risks, and one of five
outcomes: accept, provisional, reject, subsume, or no new Primitive justified.

\subsection{KFD-6: a bounded experimental loop}

At loop step $n$, let $E_n$ be an immutable causal-experience cut, $O_n$ the
declared ontology, $G_n$ plural generation experiments, $C_n$ a shared-budget
method comparison, $Q_n$ KFD-5 qualification, $H_n$ held-out or independent
evaluation, and $A_n$ promotion authority. Draft KFD-6 proposes:

\begin{equation}
(E_n,O_n) \rightarrow G_n \rightarrow C_n \rightarrow Q_n
\rightarrow H_n \rightarrow A_n \rightarrow O_n \;\text{or}\; O_{n+1}.
\end{equation}

The conservative baselines are part of the mechanism:

\begin{itemize}[leftmargin=*]
  \item \textbf{fixed ontology:} improve prediction, retrieval, or workflow
        while forbidding new first-class objects;
  \item \textbf{no new Primitive:} permit diagnosis and local abstraction but
        require the final decision to retain $O_n$; and
  \item \textbf{flattened history:} use anonymous events or endpoint facts
        under the same resource budget.
\end{itemize}

Candidates are evaluated for yield, false-candidate rate, ontology distance,
compression, intervention value, transfer, falsifiability, cost, and
auditability. Negative results remain part of the cut. Generated examples
cannot be the only evidence. The generator cannot be the sole verifier or
promotion authority.

\subsection{The authority flow}

The complete mechanism assigns a narrow authority to each layer:

\begin{center}
\begin{tabular}{p{0.20\linewidth}p{0.68\linewidth}}
\toprule
Layer & Authority \\
\midrule
Episode corpus & Evidence carrier with declared capture and stewardship; not
an ontology decision. \\
KFD-4 replay & Declared projection and comparison; not source mutation or
candidate certification. \\
Genesis method & Proposes a candidate and disconfirming test; not truth or
promotion. \\
KFD-5 qualification & Judges bounded load-bearing utility against facts and
alternatives; not universal adoption. \\
Held-out evaluator & Tests transfer and false candidates outside the discovery
cut; not value or legal authority. \\
Promotion authority & Admits or rejects ontology change inside a declared
scope; does not rewrite prior Episodes. \\
\bottomrule
\end{tabular}
\end{center}

This separation is the central anti-self-certification property. It lets an
agent contribute ontology change without turning model confidence into social
or operational authority.

\section{Problem-Definition Power and Knowledge Governance}
\label{sec:question-rights}

Knowledge systems are often evaluated by answer quality while the space of
permitted questions remains implicit. Yet a representation determines which
objects can be counted, compared, assigned, challenged, or governed. The party
that controls the durable schema can make some burdens visible and leave
others as repeated human reconstruction.

Episode infrastructure changes this structure in two opposing ways. When it
preserves multiple perspectives, causal links, responsibility, omissions, and
portable evidence, later participants can challenge the founding ontology.
When it flattens histories into provider-controlled summaries or anonymous
events, it can harden that ontology and make alternatives unverifiable. The
ability to revise problem definitions is therefore not a magical property of
memory. It is a governed capability distributed across capture, access,
representation, replay, qualification, and promotion.

Boundary-object research explains how artifacts can retain identity across
heterogeneous viewpoints without requiring consensus
\cite{star1989boundary}. KFD adds an admission question: when does a recurring
relation deserve responsibility as a Primitive, and who may decide? The answer
cannot be ``the model that named it.'' A participant has local epistemic
priority over consequences it directly bears, but its wider interpretation
remains a claim. Operators may steward shared infrastructure, but stewardship
does not authorize appropriation of private experience or unilateral ontology
change.

\subsection{Rights and obligations}

An Episode asset policy should make at least the following independently
addressable:

\begin{itemize}[leftmargin=*]
  \item the right to know whether an Episode exists and which purpose justifies
        its retention;
  \item the right to inspect, export, correct through successor evidence, or
        request deletion where applicable;
  \item the right to see redaction, degradation, provider migration, derived
        indexes, and model-generated projections;
  \item the right of a consequence-bearing participant to contest a replay or
        candidate attributed to its perspective;
  \item the obligation to preserve prior cuts rather than rewriting history to
        match a later ontology; and
  \item the separation of access to evidence, authority to qualify, and
        authority to promote.
\end{itemize}

These are design obligations, not a claim that one policy fits every legal or
cultural domain. Some Episodes should expire. Some should never be captured.
Some may support only local replay, not cross-organization training or
discovery. An honest system preserves these limits as part of the asset rather
than treating them as obstacles external to it.

\subsection{Why the issue is structurally important}

If agents participate in producing candidate concepts, ontology change becomes
part of the knowledge-production pipeline rather than a rare upstream design
event. The pipeline can redistribute who notices unnamed burdens and who can
propose new objects. It can also concentrate power in the corpus owner, model
provider, or promotion authority. That is why the problem cannot be reduced to
better retrieval. The object under design is a constitutional separation among
history, interpretation, qualification, and admission.

The appropriate conclusion is neither ``automate problem definition'' nor
``reserve concepts for humans.'' It is to let humans and agents contribute at
any stage when their evidence and perspective are declared, while preserving
independent challenge and accountable promotion for consequential changes.

\section{Bounded Evidence from Current KFD Cases}
\label{sec:evidence}

\subsection{Immutable evidence cut}

This paper freezes KFD repository commit
\texttt{ada402f671a878ff21e2f0d1c1b21c0458bac7a5}, dated July 21, 2026
\cite{kfdCut2026}. Later changes cannot be projected backward into the paper's
claim. At this cut, KFD-4 and KFD-5 are active; KFD-6 is draft. The live-case
registry contains three active cases and six candidate tracks
\cite{kfdCases2026}.

\begin{table}[ht]
\centering
\small
\begin{tabular}{p{0.265\linewidth}p{0.195\linewidth}p{0.145\linewidth}p{0.275\linewidth}}
\toprule
Live case & Track & Status & Claim boundary \\
\midrule
Proof-Carrying Work Object & Pursuit & Provisional & Durable intended-change
identity; no universal work ontology or accepted Primitive claim. \\
 & Warrant & Provisional & Purpose-bound authorization; no replacement for
enforcement or accepted Primitive claim. \\
\addlinespace
Software-Work Perspective Settlement & Initiative & Accepted & Accepted only
as a software-development Domain Profile Primitive. \\
 & Assignment & Accepted & Accepted only as bounded software-work
responsibility in the current profile. \\
 & Project Cut & Accepted & Accepted only as a software-project settlement
Primitive; not a fourth fact engine or universal object. \\
\addlinespace
Decision-Admission Settlement & Consequential Settlement & Provisional & Draft
KFD-11 procedure; no activation, universal minimality, or independent
cross-domain adoption. \\
\bottomrule
\end{tabular}
\caption{KFD live-case evidence at commit \texttt{ada402f}. ``Accepted'' is
scope-bounded and does not mean universal or historically novel.}
\label{tab:cases}
\end{table}

The table establishes public case coordinates and outcome boundaries. It does
not establish that an Episode corpus caused every candidate, that one method
was necessary or sufficient, or that acceptance predicts transfer.

\subsection{Two perspective transformations}

The software-work case records two transformations over retained development
history. Moving from Primitive developer to real human user exposed
\emph{Initiative} and \emph{Assignment}: lower action coordinates did not by
themselves state which work belonged together or whether a participant
accepted responsibility. Moving from coordinator to successor acting agent
exposed \emph{Project Cut}: available source, Episode, Atlas, policy, omission,
and risk records did not by themselves identify the one combination officially
settled for continuation.

Each accepted track includes a deletion witness. If Initiative grouping is
erased, equal lower objects can still require different coordination. If
Assignment acceptance is erased, equal objective and authority can still
require different accountability. If Project Cut selection is erased, equal
components can still distinguish a verified settlement from a candidate
combination. The qualifications remain limited to their software profiles.

\subsection{Natural pre-enactment of part of KFD-6}

The Consequential Settlement case is the strongest evidence for this paper's
specific thesis \cite{kfdConsequential2026}. During real development, agents
generated Claim, purpose-bound Assessment, and Decision before the maintainer
became aware of those named objects. Later agents replayed retained development
lineage, connected the intermediate ontology to Fact Admission, and compressed
it into a provisional procedure. Maintainer authorization, repository review,
and merge admission remained separate from generation.

This supports a bounded statement: an agent-generated ontology change can
emerge under real work pressure, remain recoverable in retained history, and
later receive separate KFD-5 qualification with less direct human translation
than the founding pattern assumed.

It is not a conforming KFD-6 experiment. The experience cut was reconstructed
after action. Methods and budgets were not preregistered. There was no plural
method comparison against fixed-ontology and no-new-Primitive baselines. The
evaluation was not held out from discovery, systematic false-candidate
rejection was absent, and no autonomous trigger began ontology review. The
case is a retrospective fixture for designing an experiment, not activation
evidence.

\subsection{Evidence grades}

We use four grades to prevent rhetorical promotion:

\begin{description}[leftmargin=2.0em,style=nextline]
  \item[E0 -- repository fact.] Exact commit, schema, case, cut digest, status,
        test, review, or release coordinate.
  \item[E1 -- implementation-backed case inference.] A bounded interpretation
        supported by retained public history and explicit falsifiers.
  \item[E2 -- prospective hypothesis.] A mechanism or metric requiring a new
        controlled or shadow evaluation.
  \item[E3 -- non-claim.] Universal validity, historical novelty, general
        autonomous discovery, or quantified social value not supported here.
\end{description}

Table~\ref{tab:cases} is E0. The natural pre-enactment interpretation is E1.
The discovery option model and prospective evaluation are E2. A statement that
KFD-6 is operationally proven would be E3 and is rejected.

\section{Prospective Evaluation Protocol}
\label{sec:evaluation}

A useful evaluation must distinguish three propositions: the Episode
representation preserves information unavailable in flattened histories; the
discovery mechanism uses that information better than conservative baselines;
and the resulting candidate transfers without exporting more complexity or
governance harm than it compresses.

\subsection{Protocol A: representation ablation}

Construct several immutable corpora from the same completed work Episodes:

\begin{enumerate}[leftmargin=*]
  \item full Episode assets with perspective, causality, responsibility,
        evidence, gaps, and stewardship metadata;
  \item provenance-preserving events without participant consequence and
        ontology declarations;
  \item anonymous event sequences with timestamps and payloads; and
  \item endpoint Fact cuts only.
\end{enumerate}

Blind evaluators to the known candidate names. Ask them to reconstruct why
specified decisions differed, identify repeated unnamed burdens, and propose
or reject candidate distinctions. Measure reconstruction time, unresolved
questions, evidence alignment, candidate yield, false-candidate rate, and
inter-evaluator agreement. If the full Episode condition does not improve
recoverability or reduces it by imposing excessive structure, RQ1 is weakened.

The ablation must hold source cases and resource budgets constant. Otherwise a
larger corpus or more capable model could be mistaken for a representation
effect.

\subsection{Protocol B: preregistered shadow KFD-6 loop}

Before new work begins, freeze $E_n$ and $O_n$, preregister method procedures,
budgets, metrics, falsifiers, and autonomy boundaries, and reserve a held-out
cut. Run at least two method families, for example perspective transformation
and anomaly/reconstruction analysis, plus fixed-ontology and no-new-Primitive
baselines.

Candidate names and attractive narratives are not the primary endpoint. The
endpoint is whether a candidate survives KFD-5 minimum closure, alternative,
deletion, fuse, falsifier, and dogfood gates and then improves a held-out
decision. Retain rejected candidates and score why they failed. A system that
generates many plausible objects but cannot reject them fails the experiment.

Shadow mode prohibits the loop from changing production ontology or participant
responsibility. An independent evaluator receives the held-out evidence. A
separate accountable authority decides whether any result warrants a bounded
intervention.

\subsection{Protocol C: transfer and intervention}

A provisional candidate that passes Protocol B is mapped into at least one
context outside its founding workstream. The test permits different names and
representations. It asks whether an equivalent distinction:

\begin{itemize}[leftmargin=*]
  \item reduces repeated reconstruction or coordination cost;
  \item changes a decision that remained ambiguous under the prior ontology;
  \item preserves source Fact, Episode, perspective, and authority boundaries;
  \item admits a smaller rival representation if one works better; and
  \item remains removable or reversible when consequences disagree with the
        qualification.
\end{itemize}

Transfer within one organization's repositories is weaker than independent
adoption. Both are useful but must not share a label.

\subsection{Protocol D: governance and stewardship}

The experiment audits the corpus itself. For sampled Episodes, verify integrity
roots, retention basis, consent or other legitimate capture boundary,
redaction, deletion behavior, export/import fidelity, provider migration, and
the distinction between authoritative records and derived projections. Invite
consequence-bearing participants to challenge replays attributed to them.

Governance failure can invalidate a technically promising result. A candidate
is not qualified if producing it requires undeclared surveillance, prevents
applicable deletion, silently broadens purpose, or makes independent challenge
impossible.

\subsection{Decision rule}

The strong hypothesis survives only if the Episode condition produces a
candidate that passes KFD-5, improves at least one held-out decision over both
conservative baselines, rejects attractive false candidates, transfers across
a declared boundary, and preserves governance constraints at acceptable cost.
Failure should narrow the representation, mechanism, or domain claim. It must
not be repaired by relabeling retrieval or clustering as Primitive discovery.

\section{Related Work}
\label{sec:related}

\paragraph{Provenance and event history.}
W3C PROV provides a portable model for entities, activities, agents,
derivation, and responsibility \cite{w3cProv2013}. Process mining discovers
and checks process models from event data and has long emphasized event-log
quality and fitness between observed and modeled behavior
\cite{vanDerAalst2012manifesto}. KFD Episodes can use both. The additional
problem is that the case, activity, responsibility, or process ontology may be
the object under revision. KFD therefore requires the current ontology,
observer, consequence boundary, and promotion path to remain explicit.

\paragraph{Causal discovery and representation learning.}
Causal inference distinguishes observation, intervention, and counterfactual
reasoning \cite{pearl2009causality}. Causal representation learning asks how
high-level causal variables can be recovered from low-level observations and
how they support transfer \cite{scholkopf2021causal}. This is close to
Primitive genesis, but KFD addresses a wider socio-technical object: a
candidate may carry responsibility, authority, lifecycle, and participant
consequences that cannot be promoted by model fit alone. KFD-5 qualification
and separated promotion authority are therefore complementary to statistical
causal evidence.

\paragraph{Episodic control and open-ended learning.}
Model-free episodic control stores high-value experience for rapid reuse
\cite{blundell2016episodic}. Open-ended learning changes task distributions and
objectives so agents continue acquiring behavior
\cite{stooke2021open,guttenberg2019open}. These lines demonstrate the value of
experience beyond one update and the difficulty of measuring continuing
novelty. The present proposal is narrower in capability and broader in
governance. It does not seek unbounded behavioral novelty; it tests whether a
declared causal corpus supports a bounded ontology change that remains
falsifiable and separately admitted.

\paragraph{Organizational memory.}
Organizational-memory research distinguishes acquisition, retention,
retrieval, use, misuse, and abuse \cite{walsh1991memory}. Episode asset
stewardship shares this concern but rejects the repository metaphor when it
implies a neutral past. Perspective and schema determine what enters memory,
while later replay constructs a view under declared loss. The asset is thus a
versioned evidence option, not a complete organizational recollection.

\paragraph{Boundary objects and distributed cognition.}
Boundary objects are adaptable to different social worlds while maintaining
enough common identity for cooperation \cite{star1989boundary}. KFD's accepted
software-domain candidates often mediate exactly such contact surfaces. The
new contribution is a promotion discipline: adaptability and use do not alone
establish a Primitive. The candidate must preserve a decision-changing
distinction, survive alternatives and falsifiers, and keep its lower
authorities inspectable.

Across these literatures, the common insight is that history, representation,
and viewpoint shape action. The gap addressed here is an end-to-end contract
for moving from governed causal history to a proposed ontology revision
without merging generation, evidence, evaluation, and admission.

\section{Limitations, Failure Modes, and Non-Claims}
\label{sec:limitations}

\subsection{Evidence is first-party and retrospective}

The current live cases are maintained by the same ecosystem that developed
KFD. Public commits, immutable cuts, review records, and explicit falsifiers
make the claims inspectable but do not remove correlated incentives. The
Consequential Settlement mapping was reconstructed after action. Independent
adopters and prospective trials are required.

\subsection{Survivorship and narrative bias}

Once a candidate is named, many histories can be retold as if they led
inevitably to it. The mechanism counters this with frozen cuts, method
plurality, conservative baselines, false-candidate retention, and held-out
evaluation, but the current evidence does not yet demonstrate that those gates
are sufficient.

\subsection{Capture can distort action}

Rich Episode capture may change participant behavior, create surveillance
pressure, and privilege what is machine-recordable. A causal graph can still
exclude consequences that occur outside the capture boundary. Declared gaps do
not eliminate this harm. In sensitive settings, less retention or no capture
may be the correct choice.

\subsection{Ontology growth can export complexity}

A false Primitive replaces local reconstruction with global schema,
governance, migration, and training cost. Deletion and fuse tests are
necessary, not sufficient. Progressive disclosure and rival representations
must remain available. ``No new Primitive'' is a successful outcome.

\subsection{No automatic epistemic justice}

Perspective metadata does not guarantee that marginalized participants can
contribute, contest, or control promotion. The corpus owner may still dominate
access and interpretation. The governance requirements in
Section~\ref{sec:question-rights} are research and design obligations, not a
proof that a deployed institution is legitimate.

\subsection{Explicit non-claims}

This alpha paper does not claim:

\begin{itemize}[leftmargin=*]
  \item that KFD-6 is active or that a conforming autonomous loop exists;
  \item that every trace should become an Episode or every Episode be retained;
  \item that more data produces more Primitives;
  \item that perspective transformation dominates other genesis methods;
  \item that the six current candidate tracks are historically novel,
        universally minimal, or cross-domain requirements;
  \item that a schema validator proves causal truth, discovery, or promotion;
  \item that one physical object model is required; or
  \item that technical stewardship resolves legal, ethical, labor, cultural,
        or political authority.
\end{itemize}

These boundaries are part of the result. A theory of ontology revision that
cannot preserve the possibility that it is wrong would reproduce the failure
it claims to solve.

\section{Conclusion}
\label{sec:conclusion}

Historical assets are usually valued for remembering what a system already
knew how to represent. This paper proposes a stronger and more demanding role:
a governed causal Episode corpus can preserve the option to discover that the
old representation itself was inadequate.

That option requires structure and restraint. KFD-4 makes replay
perspective-bound and loss-aware. KFD-5 keeps candidate genesis open while
making qualification answerable to facts, alternatives, deletion, falsifiers,
and real work. Draft KFD-6 adds plural method comparison, fixed-ontology and
no-new-Primitive baselines, held-out evaluation, false-candidate retention,
and separated promotion authority. Episode storage alone performs none of
these acts.

The current KFD case registry makes the research direction more visible. It
contains six real candidate tracks, including three bounded software-domain
acceptances and a retained natural pre-enactment of agent-generated ontology
change. It also makes the missing evidence visible: no preregistered plural
method experiment, no systematic false-candidate result, no held-out transfer,
and no autonomous ontology-review trigger. That is why the evidence supports
feasibility, not completion.

The practical consequence is significant. If agents can participate in
defining the objects through which work is understood, problem definition
becomes an operational knowledge-production surface. Preserving that surface
can expand who notices unnamed burdens; governing it poorly can concentrate
power in corpus, model, or promotion owners. The right response is neither
unbounded retention nor unaccountable automation. It is a verifiable
separation among causal history, perspective, candidate generation,
qualification, independent evaluation, and admission.

The next milestone is therefore empirical: run the preregistered ablation and
shadow-loop protocols in Section~\ref{sec:evaluation}, retain the negative
results, and let failure narrow the claim. If a candidate survives those gates
and improves held-out continuation without hiding participant or source
authority, Episode history will have demonstrated value not merely as memory,
but as infrastructure for revising what a system is capable of knowing.


\begingroup
\raggedright
\bibliographystyle{plain}
\bibliography{references}
\endgroup

\end{document}
